% ============================================================================
\section{Latar belakang dan rumusan masalah}
\label{sec:masalah}
% ============================================================================

\subsection{Perhitungan PPh 21 pegawai tetap}
\label{subsec:pph21}

PPh 21 pegawai tetap dihitung dari bruto, yaitu penghasilan teratur seperti gaji, tunjangan, dan lembur, penghasilan tidak teratur seperti THR dan bonus, serta premi jaminan yang dibayar pemberi kerja~\cite{pmk168_2023}. Neto adalah bruto dikurangi biaya jabatan dan iuran pensiun pegawai. Penghasilan kena pajak (PKP) adalah neto setahun dikurangi penghasilan tidak kena pajak~\cite{pmk101_2016}, dibulatkan ke bawah hingga ribuan rupiah, lalu dikenai tarif Pasal 17 yang terdiri atas lima lapisan dari 5\% sampai 35\%~\cite{uu7_2021}. Sampai 2023, potongan bulanan dihitung dengan menyetahunkan penghasilan~\cite{per16_2016}. Sejak 2024, potongan Januari sampai November adalah TER dikali bruto bulanan, dengan kategori A, B, dan C menurut status PTKP~\cite{pp58_2023}, dan selisihnya terhadap PPh setahun dipotong di masa pajak terakhir~\cite{pmk168_2023}. Pada 2025--2026, PPh 21 pegawai di sektor tertentu ditanggung pemerintah~\cite{pmk10_2025,pmk105_2025}. Tabel~\ref{tab:ruang-pajak} merangkum dimensi yang menentukan besarnya pajak. Penelitian ini terbatas pada pegawai tetap dengan satu pemberi kerja pada tahun pajak 2023--2026.

\begin{table}[t]
\centering
\caption{Dimensi perhitungan PPh 21 pegawai tetap yang dicakup}
\label{tab:ruang-pajak}
\scriptsize
\begin{tabularx}{\columnwidth}{@{}P{2.0cm} L@{}}
\toprule
\textbf{Dimensi} & \textbf{Nilai} \\
\midrule
Rezim & Disetahunkan (sampai 2023); TER (sejak 2024); TER dengan DTP (2025--2026) \\
Lapisan Pasal 17 & 5\%, 15\%, 25\%, 30\%, 35\% \\
Kategori TER & A, B, C dengan 44, 40, dan 41 lapisan tarif bulanan \\
Status PTKP & TK/0--TK/3, K/0--K/3; menentukan kategori TER \\
Metode & Dipotong (\en{gross}), \en{gross-up}, ditanggung pemberi kerja \\
Pola kerja & Setahun penuh, masuk atau keluar di tengah tahun \\
Penghasilan & Teratur; tidak teratur (THR, bonus) \\
\bottomrule
\end{tabularx}
\end{table}

\subsection{Kebijakan perusahaan dan batas kewenangannya}
\label{subsec:kebijakan}

Perusahaan menetapkan komponen penghasilan, seperti prorata gaji, THR, dan tunjangan, lewat peraturan perusahaan atau perjanjian kerja bersama, sehingga kebijakan itu ikut menentukan PPh 21. Namun, perusahaan tidak boleh menyimpangi ketentuan wajib. Contohnya, PER-16/PJ/2016 dan PMK 168/2023 menggolongkan uang lembur sebagai penghasilan teratur~\cite{per16_2016,pmk168_2023}, sedangkan kebijakan payroll perusahaan yang diaudit mencatatnya sebagai penghasilan tidak teratur. Bila kebijakan bertentangan dengan ketentuan wajib, ketentuan regulasi yang berlaku (\en{lex superior}), dan konfliknya harus tercatat.

\subsection{Jenis perubahan dan rumusan masalah}
\label{subsec:rumusan}

Perubahan dikelompokkan menjadi perubahan parameter, tabel tarif, klasifikasi penghasilan, aturan atau fasilitas baru, kondisi atau cakupan, pergantian metode, pencabutan, ketentuan berlaku surut, dan perubahan kebijakan perusahaan. Satu peristiwa dapat termasuk beberapa jenis, misalnya PMK 10/2025 yang menambah fasilitas baru sekaligus berlaku surut. Hasil dilaporkan per jenis, sehingga perubahan parameter yang mudah tidak digabung dengan pergantian metode.

Masukan masalah ini adalah \en{knowledge base} $\K_{n-1}$ yang benar untuk keadaan hukum sebelum peristiwa ke-$n$, beserta dokumen sumber peristiwa itu. Keluarannya adalah $\K_n$ yang benar untuk keadaan hukum sesudahnya, tanpa perubahan kode \en{engine}, tanpa perubahan hasil untuk masa sebelum perubahan berlaku, dan dengan laporan kebijakan perusahaan yang menjadi bertentangan dengan ketentuan wajib. Masalah ini berulang, karena $\K_n$ menjadi masukan untuk peristiwa berikutnya, sehingga kesalahan yang lolos dapat merambat.
